% Generated draft: requires mathematical and source review.
\documentclass[11pt]{article}
\usepackage[margin=25mm]{geometry}
\usepackage{fontspec}
\setmainfont{Latin Modern Roman}
\usepackage{amsmath,amssymb,mathtools,unicode-math}
\setmathfont{Latin Modern Math}
\usepackage{enumitem,needspace,xurl,hyperref}
\setlength{\parindent}{0pt}
\setlength{\parskip}{5pt}
\setlength{\emergencystretch}{3em}
\newcommand{\N}{\mathbb{N}}
\newcommand{\Z}{\mathbb{Z}}
\newcommand{\Q}{\mathbb{Q}}
\newcommand{\R}{\mathbb{R}}
\newcommand{\C}{\mathbb{C}}
\newcommand{\F}{\mathbb{F}}
\newcommand{\A}{\mathbb{A}}
\newcommand{\PP}{\mathbb{P}}
\newcommand{\E}{\mathbb{E}}
\newcommand{\eps}{\varepsilon}
\newcommand{\dd}{\,\mathrm{d}}
\begin{document}
\section*{UnsolvedMath ID: 30003071}
\textbf{Source problem code:} catalogue entry
\subsection{Definitions and mathematical statement}
Write \(\Gamma(X,\mathcal F)\) for the space of global sections of a sheaf \(\mathcal F\) on \(X\), and \(\operatorname{Sym}^{\bullet}\mathcal E=\bigoplus_{d\geq 0}\operatorname{Sym}^{d}\mathcal E\) for the graded symmetric algebra of a specified sheaf \(\mathcal E\). The notation \(GL_N\) denotes a general linear group, but the source record does not specify its base field or the admissible values of \(N\) and \(r\).

The supplied question asks whether the conjectural Kostka--Shoji formula yields the following graded \(GL_N\)-module identity:
\[
\begin{aligned}
\Gamma\bigl(T_r^*B_N^r,O(\mu)\bigr)
&=\Gamma\bigl(B_N^r,\operatorname{Sym}^{\bullet}T_rB_N^r\otimes O(\mu)\bigr)\\
&=\bigoplus_{\lambda\geq\mu}K_{\lambda\mu}^{-,\bullet}\otimes V_\lambda.
\end{aligned}
\]
Here the symbols \(B_N^r\), \(T_rB_N^r\), \(T_r^*B_N^r\), \(O(\mu)\), \(K_{\lambda\mu}^{-,\bullet}\), and \(V_\lambda\) are retained exactly as source notation. Their constructions, the allowed indices \(\lambda,\mu\), the relation \(\lambda\geq\mu\), the tensor product's base field, and the group actions and gradings are not defined in the supplied record. In particular, the record does not specify whether \(K_{\lambda\mu}^{-,\bullet}\) denotes a graded multiplicity space or a polynomial encoding one, or how \(O(\mu)\) on the first space relates to its namesake on the second. The displayed equalities reproduce the source's proposed identity; whether they mean canonical identifications, graded module isomorphisms, or character identities requires clarification.

No precise universal quantification or complete conjectural hypothesis can be supplied from this record alone. The background further describes the extracted passage as a corollary fragment and discusses finite generation of a monoid associated with a birational sequence, without specifying that monoid or its connection to this identity. These unresolved qualifications prevent a standalone precise formulation.
\subsection{Short English statement}
The supplied question asks whether a conjectural Kostka--Shoji formula gives the displayed graded \(GL_N\)-module decomposition of global sections. Essential objects, parameter ranges, hypotheses, and grading conventions are unspecified, and the background discusses a different, unexplained finite-generation target. The formula is therefore retained as a qualified source question requiring mathematical review.
\subsection{Sources}
\begin{itemize}
\item Supplied local UnsolvedMath registry record, permanent ID \(30003071\), source problem code \texttt{catalogue entry}, title ``Kostka--Shoji Formula for Graded Modules,'' in the Oberwolfach Reports Open Problems collection. The record gives proposed year 2016 and names no proposer.
\item Supplied report citation: \emph{Mini-Workshop: PBW Structures in Representation Theory} (2016), Oberwolfach Reports, pp. 615--651, DOI 10.4171/owr/2016/13. This is also the literature-source entry labelled ``OWR report 13/2016'': \url{https://doi.org/10.4171/owr/2016/13}. No exact passage or individual contribution is supplied.
\item Supplied literature-source entry labelled ``Fang--Fourier--Littelmann'': \url{https://arxiv.org/abs/1510.02295}. The record supplies no title or passage connecting this reference to the displayed identity.
\item Literature assessment checked at \texttt{2026-08-21}: this is supplied dated curation, not independently verified. It labels the problem open, describes the extracted passage as a corollary fragment rather than a standalone assertion, and reports that an underlying finite-generation claim for an associated birational-sequence monoid is not known in general. The supplied background identifies the report and Fang--Fourier--Littelmann entry as sources checked and requests independent verification of time-sensitive status and cited claims.
\item The registry also carries the supplied statement-status label \texttt{corrected\_verified} and a note claiming that manual semantic, typographic, and LaTeX review confirmed self-containment and preservation of the mathematical target. That supplied verification claim has not been independently confirmed and does not resolve the missing definitions or the discrepancy with the background assessment.
\end{itemize}
\end{document}
